% Generated draft: requires mathematical and source review.
\documentclass[11pt]{article}
\usepackage[margin=25mm]{geometry}
\usepackage{fontspec}
\setmainfont{Latin Modern Roman}
\usepackage{amsmath,amssymb,mathtools,unicode-math}
\setmathfont{Latin Modern Math}
\usepackage{enumitem,needspace,xurl,hyperref}
\setlength{\parindent}{0pt}
\setlength{\parskip}{5pt}
\setlength{\emergencystretch}{3em}
\newcommand{\N}{\mathbb{N}}
\newcommand{\Z}{\mathbb{Z}}
\newcommand{\Q}{\mathbb{Q}}
\newcommand{\R}{\mathbb{R}}
\newcommand{\C}{\mathbb{C}}
\newcommand{\F}{\mathbb{F}}
\newcommand{\A}{\mathbb{A}}
\newcommand{\PP}{\mathbb{P}}
\newcommand{\E}{\mathbb{E}}
\newcommand{\eps}{\varepsilon}
\newcommand{\dd}{\,\mathrm{d}}
\begin{document}
\section*{UnsolvedMath ID: 30001083}
\textbf{Source problem code:} catalogue entry
\subsection{Definitions and mathematical statement}
For a tessellation of $\R^d$, its zero cell is the cell containing the origin, unique almost surely in the usual stationary models. It is not the Palm typical cell. For a bounded $d$-polytope $P$, let $f_\ell(P)$ denote the number of its $\ell$-dimensional faces. The source poses two separate moment questions.\begin{enumerate}\item For a Poisson hyperplane tessellation, write $Z_0^{\mathrm H}$ for its zero cell. In the standard construction it is the intersection of the halfspaces containing the origin bounded by the process hyperplanes. Determine $\E f_\ell(Z_0^{\mathrm H})$ for $\ell\geq1$. The nontrivial proper-face range is $1\leq\ell\leq d-1$; if the whole polytope is counted as a face, $f_d=1$, and $f_\ell=0$ for $\ell>d$. The dimension, hyperplane intensity/directional distribution, and nondegeneracy assumptions making the cell bounded are not given. No isotropic specialization is silently imposed.\item For a Poisson--Voronoi mosaic, write $Z_0^{\mathrm V}$ for its zero cell and determine its mean volume and face numbers. In the usual Euclidean model generated by a homogeneous Poisson point process $\xi$ of intensity $\lambda>0$, cells are \[C(x)=\{y\in\R^d:|y-x|\leq|y-z|\text{ for every }z\in\xi\},\qquad x\in\xi,\] and $Z_0^{\mathrm V}$ is the cell containing the origin. The requested quantities are $\E\operatorname{Vol}_d(Z_0^{\mathrm V})$ and $\E f_\ell(Z_0^{\mathrm V})$ in the intended face range. This sampling differs from adjoining a nucleus at the origin and taking its cell; typical-cell moments are not substituted. The source does not state the dimension or intensity normalization or explicitly specify stationarity.\end{enumerate} The questions seek exact mean quantities for the intended models, not only large-cell asymptotics. The supplied 2026-08-21 triage labels the compound record partially solved, reports a proved large Poisson--Voronoi circle-shape conjecture and later large-zero-cell asymptotics, and says the extracted passage also recorded missing exact zero-cell moments. That separate shape/asymptotic progress is preserved as a qualification, not treated as an answer to every moment request or as a new displayed conjecture.
\subsection{Short English statement}
Determine mean nonvertex face numbers for Poisson-hyperplane zero cells and mean volume and face numbers for Poisson--Voronoi zero cells. The dimension and process parameters are unstated, and zero-cell sampling must not be replaced by typical-cell sampling.
\subsection{Sources}
\begin{itemize}\item Local UnsolvedMath registry, permanent ID 30001083, source problem catalogue entry, \emph{Face Numbers and Volumes of Poisson Zero Cells}, proposed year 2008; no individual proposer supplied.\item Supplied citation: \emph{New Perspectives in Stochastic Geometry} (2009), Oberwolfach Reports, DOI 10.4171/owr/2008/47, pp.\ 2655--2702: \url{https://doi.org/10.4171/owr/2008/47}. Supplied literature label: original OWR report.\item Supplied literature label: Hug--Reitzner--Schneider, large Poisson--Voronoi cells, \url{https://doi.org/10.1239/aap/1093962228}.\item Supplied literature label: Calka--Demichel--Enriquez, large planar zero cells, \url{https://www.numdam.org/item/10.5802/ahl.86.pdf}.\item Literature assessment date: 2026-08-21. The status and progress claims are supplied dated curation, not independently verified here.\end{itemize}
\end{document}
